%%%%%%%%%%%%%%%%%%%%%%%%%%%%%%%%%%%%%%%%%%%%%%%%%%%%%%%%%%%%%%%%%%%%%%%%%%% 
% 
% Generic template for the anteproyectos of TFC/TFM/TFGs
% 
% $Id: anteproyecto.tex,v 1.6 2018/09/11 12:23:48 macias Exp $
% 
% By:
%  + Javier Macías-Guarasa. 
%    Departamento de Electrónica
%    Universidad de Alcalá
%  + Roberto Barra-Chicote. 
%    Departamento de Ingeniería Electrónica
%    Universidad Politécnica de Madrid   
% 
% Based on original sources by Roberto Barra, Manuel Ocaña, Jesús Nuevo,
% Pedro Revenga, Fernando Herránz and Noelia Hernández. Thanks a lot to
% all of them, and to the many anonymous contributors found (thanks to
% google) that provided help in setting all this up.
% 
% See also the additionalContributors.txt file to check the name of
% additional contributors to this work.
% 
% If you think you can add pieces of relevant/useful examples,
% improvements, please contact us at (macias@depeca.uah.es)
% 
% You can freely use this template and please contribute with
% comments or suggestions!!!
% 
%%%%%%%%%%%%%%%%%%%%%%%%%%%%%%%%%%%%%%%%%%%%%%%%%%%%%%%%%%%%%%%%%%%%%%%%%%% 

% This is for rubber to clean additional files
% rubber: clean anteproyecto.acn anteproyecto.acr anteproyecto.alg anteproyecto.cod anteproyecto.ist anteproyecto.out anteproyecto.sbl anteproyecto.slg anteproyecto.sym anteproyecto.lor

\documentclass[11pt,a4paper,oneside]{article}

%%%%%%%%%%%%%%%%%%%%%%%%%%%%%%%%%%%%%%%%%%%%%%%%%%%%%%%%%%%%%%%%%%%%%%%%%%% 
% BEGIN Preamble and configuration section
% 
\input{../Config/preamble.tex}    % DO NOT TOUCH THIS LINE. You can edit
                                  % the file to modify some default settings

\input{../Config/myconfig.tex}    % DO NOT TOUCH THIS LINE, but EDIT THIS FILE
                                  % to set your specific settings (related
                                  % to the document language, your degree,
                                  % document details (such as title, author
                                  % (you), your email, name of the tribunal
                                  % members, document year, keyword and
                                  % palabras clave) and link colors), and
                                  % define your commonly used commands
                                  % (some examples are provided).

\input{../Config/postamble.tex}   % DO NOT TOUCH THIS LINE. Yes, I know,
                                  % "postamble" is not a valid word... :-)
%
% END Preamble and configuration section
%%%%%%%%%%%%%%%%%%%%%%%%%%%%%%%%%%%%%%%%%%%%%%%%%%%%%%%%%%%%%%%%%%%%%%%%%%%

% path to directories containing images
\graphicspath{{../Book/logos/}{../Book/figures/}{../Book/diagrams/}} % Edit this to your
                                  % needs. Only logos is really required
                                  % when you generate your own content.
% 
% END Preamble and configuration section
%%%%%%%%%%%%%%%%%%%%%%%%%%%%%%%%%%%%%%%%%%%%%%%%%%%%%%%%%%%%%%%%%%%%%%%%%%% 

\title{Anteproyecto de \myWorkTypeFull}        % DO NOT TOUCH THIS LINE
\date{\myThesisProposalDate}                         % DO NOT TOUCH THIS LINE
\author{\myAuthorFullName}

%%%%%%%%%%%%%%%%%%%%%%%%%%%%%%%%%%%%%%%%%%%%%%%%%%%%%%%%%%%%%%%%%%%%%%%%%%% 
% Let's start with the real stuff
%%%%%%%%%%%%%%%%%%%%%%%%%%%%%%%%%%%%%%%%%%%%%%%%%%%%%%%%%%%%%%%%%%%%%%%%%%% 
\begin{document}                                   % DO NOT TOUCH THIS LINE

%%%%%%%%%%%%%%%%%%%%%%%%%%%%%%%%%%%%%%%%%%%%%%%%%%%%%%%%%%%%%%%%%%%%%%%%%%% 
% BEGIN within-document configuration, frontpage and cover pages generation
% 

% Set Language dependent issues that must be set after \begin{document}
\input{../Config/setlanguagedependentissues.tex} % DO NOT TOUCH THIS LINE
                                                 % NOR THE FILE

% 
% END within-document configuration, frontpage and cover pages generation
%%%%%%%%%%%%%%%%%%%%%%%%%%%%%%%%%%%%%%%%%%%%%%%%%%%%%%%%%%%%%%%%%%%%%%%%%%% 

\maketitle

\begin{description}                               % DO NOT TOUCH THIS LINE
\item[\expandafter\makefirstuc\expandafter{\wordAutorOrAutora}:] \myAuthorFullName                       % DO NOT TOUCH THIS LINE
\item[\expandafter\makefirstuc\expandafter{\wordTutorOrTutores}:] \myAdvisors                   % DO NOT TOUCH THIS LINE
  \item[Titulación:] \myDegreefull                % DO NOT TOUCH THIS LINE
  \item[Título:] \myBookTitleSpanish              % DO NOT TOUCH THIS LINE
  \ifthenelse{\equal{\myLanguage}{english}}       % DO NOT TOUCH THIS LINE
  {                                               % DO NOT TOUCH THIS LINE
  \item[Título en inglés:] \myBookTitleEnglish    % DO NOT TOUCH THIS LINE
  }                                               % DO NOT TOUCH THIS LINE
  {                                               % DO NOT TOUCH THIS LINE
  }                                               % DO NOT TOUCH THIS LINE
\item[Departamento:] \myDepartment                % DO NOT TOUCH THIS LINE
\end{description}                                 % DO NOT TOUCH THIS LINE

%\tableofcontents

%%%%%%%%%%%%%%%%%%%%%%%%%%%%%%%%%%%%%%%%%%%%%%%%%%%%%%%%%%%%%%%%%%%%%%%%%%% 
%%%%%%%%%%%%%%%%%%%%%%%%%%%%%%%%%%%%%%%%%%%%%%%%%%%%%%%%%%%%%%%%%%%%%%%%%%% 
%%%%%%%%%%%%%%%%%%%%%%%%%%%%%%%%%%%%%%%%%%%%%%%%%%%%%%%%%%%%%%%%%%%%%%%%%%% 
%%%%%%%%%%%%%%%%%%%%%%%%%%%%%%%%%%%%%%%%%%%%%%%%%%%%%%%%%%%%%%%%%%%%%%%%%%% 
%%%%%%%%%%%%%%%%%%%%%%%%%%%%%%%%%%%%%%%%%%%%%%%%%%%%%%%%%%%%%%%%%%%%%%%%%%% 
%%%%%%%%%%%%%%%%%%%%%%%%%%%%%%%%%%%%%%%%%%%%%%%%%%%%%%%%%%%%%%%%%%%%%%%%%%% 
%%%%%%%%%%%%%%%%%%%%%%%%%%%%%%%%%%%%%%%%%%%%%%%%%%%%%%%%%%%%%%%%%%%%%%%%%%% 
% BEGIN Normal sections. Edit/modify all within this section

\section{Introducción}
\label{sec:introduccion}

En los últimos años, la Inteligencia Artificial (IA) ha dejado de ser un concepto puramente académico para convertirse
en un pilar fundamental de la transformación tecnológica en la actualidad. Aunque para la mayoría, la IA se asocia
inmediatamente con modelos generativos de texto como ChatGPT o Gemini, o imágenes llamativas que acaparan la atención mediática,
pero su verdadero impacto en la ingeniería radica en tareas como como procesamiento de grandes volúmenes de datos, la  
automatización de procesos complejos, o la toma de decisiones óptimas. Dentro de este entorno, existen diferentes
vertientes como el aprendizaje supervisado, el no supervisado, o el aprendizaje por refuerzo, las cuales ofrecen
herramientas muy versátiles que permiten resolver problemas analíticos de diversa naturaleza de una manera mucho más 
dinámica que la programación tradicional.

Uno de los campos donde la IA está demostrando cada vez más su excepcional valor estratégico es en la optimización
combinatoria y la logística, las cuales son áreas fuertemente representadas por retos matemáticos clásicos como el Problema
del Viajante (en inglés, TSP, \textit{Traveling Salesman Problem}). Tradicionalmente, la resolución de estos problemas
de rutas ha dependido de algoritmos matemáticos exactos, y técnicas metaheurísticas rígidas, que requieren de una 
configuración manual muy específica. Sin embargo, la integración de técnicas modernas de \textit{machine learning} permite
abordar la optimización de este problema desde una nueva perspectiva, más híbrida e inteligente. Esto abre la puerta a 
utilizar el aprendizaje automático para simplificar la complejidad de los mapas de datos, o para guiar la búsqueda de caminos 
óptimos, permitiendo así la obtención de soluciones de alta calidad, sin la necesidad de depender de infraestructuras de 
\textit{hardware} masivas.

La motivación de este proyecto surge, precisamente, de la necesidad de analizar críticamente el rendimiento de estas 
diferentes soluciones en entornos de computación convencionales (CPU), centrando el esfuerzo en la replicación y 
comparación analítica de diversos algoritmos heurísticos y metaheurísticos consolidados en la literatura (tales como algoritmos 
genéticos, recocido simulado o colonia de hormigas, entre otros). Bajo esta premisa, el presente trabajo plantea el 
desarrollo de un estudio comparativo riguroso que evalúe variables clave como el tiempo de cómputo, la precisión y la tasa 
de convergencia de estos métodos al enfrentarse a diferentes instancias del problema del viajante. El objetivo principal 
es construir un marco de trabajo (\textit{framework}) ordenado y reproducible, que permita validar empíricamente los resultados 
de la literatura científica actual, permitiendo además la posibilidad de explorar como técnicas de \textit{machine learning} 
podrían integrarse en el futuro, para asistir a estos algoritmos en la toma de decisiones.

\textit{En este apartado se describirá el contexto en el que se desenvolverá el
  TFG y sus antecedentes si existen. Recuerda que en el anteproyecto
  también se pueden poner referencias bibliográficas como \cite{moore2002}}.


\section{Objetivos}
\label{sec:objetivos-y-campo}

Para abordar este trabajo, combinando la investigación teórica con su posterior desarrollo práctico, se plantean los siguientes 
objetivos:

\begin{itemize}
  \item Investigación, aprendizaje y análisis comparativo de los fundamentos teóricos de diferentes algoritmos heurísticos y 
  metaheurísticos consolidados en la optimización combinatoria, comprendiendo su comportamiento y configuración matemática frente 
  al problema del viajante.
  \item Con los aspectos teóricos estudiados, y aplicando el lenguaje de programación Python, llevarlos a la práctica mediante el 
  desarrollo de un marco de trabajo (\textit{framework}) ordenado y modular. Este sistema servirá como prototipo para evaluar el 
  rendimiento de dichos algoritmos, utilizando diferentes instancias de datos estandarizadas. Se incluyen las siguientes etapas:
  \begin{itemize}
    \item Automatización de la lectura, procesamiento y carga de los conjuntos de datos de prueba del problema.
    \item Implementación aislada e independiente de cada una de las metaheurísticas seleccionadas, para asegurar la flexibilidad 
    del sistema.
    \item Monitorización y registro preciso de métricas de rendimiento ejecutadas en entorno local, midiendo limpiamente el tiempo 
    exacto de cómputo, el coste de la ruta obtenida, y la velocidad de convergencia de cada algoritmo.
  \end{itemize}
  \item Evaluación crítica del rendimiento de las soluciones desarrolladas frente a los resultados documentados en la literatura 
  científica actual, con el fin de determinar la viabilidad y escalabilidad del prototipo, dejando planteada una pequeña base para 
  integrar, de forma complementaria, técnicas ligeras de \textit{machine learning}.
\end{itemize}

\section{Descripción del trabajo}
\label{sec:descr-del-trab}

Este trabajo tendrá como objetivo final el desarrollo de un marco de trabajo (\textit{framework}) modular para la réplica y evaluación 
analítica de algoritmos metaheurísticos aplicados al Problema del Viajante. Para lograr esto, las etapas del proyecto se estructuran de 
la siguiente manera: 

\begin{enumerate}
  \item \textbf{Estudio de algoritmos heurísticos y metaheurísticos}: consiste en la formación, estudio y comprensión de los fundamentos
  de la optimización combinatoria. Se analizará a fondo el fundamento teórico de los algoritmos seleccionados, para asentar la bases 
  antes de pasar a la fase práctica.
  \item \textbf{Análisis de herramientas y librerías de desarrollo}: se evaluará el ecosistema del lenguaje Python para identificar las
   librerías más eficientes que faciliten la estructuración del sistema, el procesamiento de datos, la medición precisa de tiempos de
    computación, y la generación automática de gráficos comparativos.
  \item \textbf{Selección y preparación de instancias de datos}: en lugar de crear un conjunto de datos desde cero, se seleccionarán y
  descargarán problemas de prueba estandarizados, procedentes de la biblioteca pública TSPlib. Se desarrollará un módulo dedicado para
  automatizar la lectura, limpieza y carga de estas instancias, garantizando así que todos los algoritmos compitan bajo las mismas 
  condiciones.
  \item \textbf{Desarrollo del \textit{framework} y pruebas iniciales}: se procederá a la implementación limpia y aislada de cada 
  algoritmo dentro del sistema. En esta fase, se realizarán ejecuciones preliminares en entorno local, para verificar que el código
  funciona correctamente, así como para calibrar los parámetros básicos de cada metaheurística y asegurar la estabilidad del
  \textit{software}.
  \item \textbf{Análisis comparativo de rendimiento}: con el sistema validado, se ejecutarán las pruebas a gran escala de forma 
  controlada. Se recopilarán, de manera automática, las métricas clave en cada ejecución para realizar una comparación estadística 
  rigurosa de los resultados.
  \item \textbf{Documentación}: implica mantener el código fuente comentado de forma clara bajo buenas prácticas de programación, junto  
  con la redacción final de la memoria del proyecto, donde se detallarán los fundamentos teóricos, la arquitectura del \textit{software}
  desarrollado, y las conclusiones del análisis comparativo.
\end{enumerate}

\textit{Aquí se explicará en detalle qué se va a hacer, cómo y por qué. Sería interesante incluir un esquema/diagrama de bloques, si se puede (como el mostrado en la figura~\ref{fig:arquitectura}).}


\begin{figure}[tphb]
  \centering
  \includegraphics[width=3in]{Diagrama2.pdf}
  \caption{Arquitectura del sistema completo.}
  \label{fig:arquitectura}
\end{figure}


\section{Metodología y plan de trabajo}
\label{sec:metodologia-y-plan}

\textit{Aquí se incluirá una descripción (puede ser incluso una enumeración)
  clara de las etapas que se van a seguir, y si es posible se deberá
  incluir un diagrama de Gantt. Por ejemplo algo del estilo a:}

Estas son las fases de desarrollo que se van a seguir para la
consecución de los objetivos del proyecto descritos en la sección~\ref{sec:objetivos-y-campo}:

\begin{enumerate}
  
\item Formación inicial (1 mes)
  
  \begin{itemize}
  \item Formación en \ldots
  \item Consulta de la API \ldots
  \item Consulta bibliográfica \ldots
  \item Profundización en herramientas de soporte \ldots
  \item  \ldots
  \end{itemize}

\item Diseño del entorno (0,5 meses)

\item Diseño, implementación y evaluación del módulo 1 (2 meses)
  \begin{itemize}
  \item Definición del  \ldots
  \item Implementación del  \ldots
  \item Evaluación del  \ldots
  \end{itemize}
  
\item Diseño, implementación y evaluación del módulo 2 (1 mes):
  \begin{itemize}
  \item Definición de  \ldots
  \item Definición de  \ldots
  \item Implementación y evaluación  \ldots
  \end{itemize}

\item Integración y pruebas (1 mes)

\item Documentación (a lo largo de todo el proyecto)

\end{enumerate}

El diagrama de Gantt asociado a esa definición es el mostrado en la figura~\ref{fig:diagrama-gantt} (\textit{para lo que he usado la funcionalidad del paquete} \texttt{pgfgantt} \textit{pero que puedes hacer con cualquier otra herramienta}).

\begin{figure}[H]
  \centering

  \begin{ganttchart}[
    title height=1,
    x unit=0.08cm, % Adjusts the width of the chart
    y unit title=0.6cm, % Adjusts the height of titles
    y unit chart=0.5cm, % Adjusts the height of tasks
    group height=0.5,
    bar height=0.5,
    vgrid={*{6}{draw=none}, *1{draw, dotted}}, % Vertical lines only at month start
    hgrid,
    time slot format=isodate,
    group label font=\bfseries\footnotesize,
    bar label font=\footnotesize,
    title label font=\footnotesize,
    group/.append style={fill=gray!75}, % Change group bar fill to gray
    group peaks height=0.3, % Increase the height of the vertical marks
    group peaks width=2, % Increase the width of the vertical marks
    group peaks tip position=0, % Center the vertical marks
    ]{2026-06-24}{2026-09-08} % Adjust dates as needed

    % Define months
    \gantttitlecalendar{year, month} \\

    % Tasks
    \ganttbar{1. Formación inicial}{2026-07-01}{2026-07-14} \\
    \ganttbar{2. Diseño de la arquitectura}{2026-07-15}{2026-07-21} \\
    \ganttbar{3. Módulo de optimización}{2026-07-22}{2026-08-04} \\
    \ganttbar{4. Módulo de benchmarking}{2026-08-05}{2026-08-19} \\
    \ganttbar{5. Integración y análisis}{2026-08-20}{2026-09-01} \\
    \ganttbar{6. Documentación}{2026-07-01}{2026-09-01}
  \end{ganttchart}
  \caption{Diagrama de Gantt de la planificación del trabajo.}
  \label{fig:diagrama-gantt}

\end{figure}

\section{Medios}
\label{sec:medios}

\textit{Aquí se describen de los medios necesarios para realizar el TFG. Por
  ejemplo:}

Las herramientas que van a ser necesarias para desarrollar este proyecto
son las siguientes:

\begin{itemize}
\item PC compatible
\item Sensor  \ldots
\item Sistema operativo GNU/Linux~\cite{gnulinux}
\item Entorno de desarrollo Emacs/Vim~\cite{emacs}
\item Procesador de textos \LaTeX~\cite{lamport94}
\item Control de versiones CVS~\cite{cvs}
\item Compilador C/C++ gcc~\cite{gcc}
\item Gestor de compilaciones make~\cite{make}
\item Robot con movilidad.
\item  \ldots
\end{itemize}



Otros recursos necesarios para la elaboración del proyecto son:

\begin{itemize}
\item Herramientas  \ldots
\item Sistema de desarrollo  \ldots
\item  \ldots
\end{itemize}




% 
% END Normal sections. Edit/modify all within this section
%%%%%%%%%%%%%%%%%%%%%%%%%%%%%%%%%%%%%%%%%%%%%%%%%%%%%%%%%%%%%%%%%%%%%%%%%%% 
%%%%%%%%%%%%%%%%%%%%%%%%%%%%%%%%%%%%%%%%%%%%%%%%%%%%%%%%%%%%%%%%%%%%%%%%%%% 
%%%%%%%%%%%%%%%%%%%%%%%%%%%%%%%%%%%%%%%%%%%%%%%%%%%%%%%%%%%%%%%%%%%%%%%%%%% 
%%%%%%%%%%%%%%%%%%%%%%%%%%%%%%%%%%%%%%%%%%%%%%%%%%%%%%%%%%%%%%%%%%%%%%%%%%% 
%%%%%%%%%%%%%%%%%%%%%%%%%%%%%%%%%%%%%%%%%%%%%%%%%%%%%%%%%%%%%%%%%%%%%%%%%%% 
%%%%%%%%%%%%%%%%%%%%%%%%%%%%%%%%%%%%%%%%%%%%%%%%%%%%%%%%%%%%%%%%%%%%%%%%%%% 
%%%%%%%%%%%%%%%%%%%%%%%%%%%%%%%%%%%%%%%%%%%%%%%%%%%%%%%%%%%%%%%%%%%%%%%%%%% 


%%%%%%%%%%%%%%%%%%%%%%%%%%%%%%%%%%%%%%%%%%%%%%%%%%%%%%%%%%%%%%%%%%%%%%%%%%% 
% Bibliography
%%%%%%%%%%%%%%%%%%%%%%%%%%%%%%%%%%%%%%%%%%%%%%%%%%%%%%%%%%%%%%%%%%%%%%%%%%% 
\input{../Book/biblio/bibliography.tex}               % EDIT this file if required



\end{document}

